\section*{Artifacts}

The repository~\cite{github_repository} provides a complete experiment reproducibility kit. The \texttt{README.md} file describes the repository contents and explains how it is organized. The kit mainly allows running the crawlers \textsf{\small SB-ORACLE}, \textsf{\small SB-CLASSIFIER}, \textsf{\small FOCUSED}, \textsf{\small TP-OFF} (sometimes referred to as \verb|dom_off|), \textsf{\small BFS}, \textsf{\small DFS}, and \textsf{\small RANDOM}. Three execution modes are available:
\begin{itemize}
  \item \emph{Local crawling}: used when a website has already been fully replicated in a local database (see Sec.~\ref{sec:exp-crawl-modalities}).
  \item \emph{Online-to-local crawling}: creates a local copy of a website using a naive crawler.
  \item \emph{Semi-online crawling}: first checks the local database and fetches the resource only if it is not present (applied to non-fully crawled websites, also described in Sec.~\ref{sec:exp-crawl-modalities}).
\end{itemize}

The folder \texttt{url\_classifier\_confusion\_matrices} contains all confusions matrices used to compute the ``MR'' metric in Sec.~\ref{subsubsection:metaparameters_result} for the URL classifier models.
The repository also contains code to reproduce the plots, information about the websites used in our experiments, and an extended version of this paper. The extended version includes: the proof of Prop.~\ref{prop:np-complete}; additional details about the websites' characteristics; the MIME type list defining the targets used in the experiments; the initial keyword set for \textsf{\small TRES}; and a complete blocklist of URL extensions and MIME types. It also provides the plots of the 8 websites that could not be included in Figure~\ref{fig:10_sites_plots}; exhaustive plots from the hyper-parameter studies; extra examples of typical tag paths for 6 websites; plots and confusion matrices for the URL classifier models evaluation; additional results on URL classification quality; a visualization of the early-stopping mechanism for two websites; a table summarizing the main characteristics of related focused crawlers; and an extended discussion of other crawler types, as well as alternatives to AUER for MAB algorithms.

A second repository~\cite{tres_modified_repository} contains a fork of the \textsf{\small TRES}~\cite{kontogiannis2022tree} crawler code, adapted for the target retrieval task described in Sec.~\ref{subsection:baselines}. 